% Generated draft: requires mathematical and source review.
\documentclass[11pt]{article}
\usepackage[margin=25mm]{geometry}
\usepackage{fontspec}
\setmainfont{Latin Modern Roman}
\usepackage{amsmath,amssymb,mathtools,unicode-math}
\setmathfont{Latin Modern Math}
\usepackage{enumitem,needspace,xurl,hyperref}
\setlength{\parindent}{0pt}
\setlength{\parskip}{5pt}
\setlength{\emergencystretch}{3em}
\newcommand{\N}{\mathbb{N}}
\newcommand{\Z}{\mathbb{Z}}
\newcommand{\Q}{\mathbb{Q}}
\newcommand{\R}{\mathbb{R}}
\newcommand{\C}{\mathbb{C}}
\newcommand{\F}{\mathbb{F}}
\newcommand{\A}{\mathbb{A}}
\newcommand{\PP}{\mathbb{P}}
\newcommand{\E}{\mathbb{E}}
\newcommand{\eps}{\varepsilon}
\newcommand{\dd}{\,\mathrm{d}}
\begin{document}
\section*{UnsolvedMath ID: 30003806}
\textbf{Source problem code:} catalogue entry
\subsection{Definitions and mathematical statement}
For a positive integer \(N\), the quantum permutation group \(S_N^+\) has coordinate \(C^*\)-algebra generated by projections \(u_{ij}\), \(1\leq i,j\leq N\), with mutually orthogonal entries in each row and column and with every row and column summing to \(1\). Its coproduct satisfies
\[
\Delta(u_{ij})=\sum_{k=1}^N u_{ik}\otimes u_{kj}.
\]
The supplied statement also uses coefficients \(u_{ij}^{n}\). Their representation labels and index ranges are not defined in the record; in particular, the superscript must not silently be interpreted as an ordinary power.

Define polynomials \(T_k\) by
\[
T_0(X)=1,\qquad T_1(X)=X,\qquad
XT_k(X)=T_{k+1}(X)+T_{k-1}(X)\quad(k\geq1).
\]
The source prescribes a random-walk state \(\varphi\) on \(S_N^+\) through
\[
\varphi(u_{ij}^{n})=\delta_{ij}\left(
\frac{N-1}{N}\frac{T_{2n}(\sqrt{N-2})}{T_{2n}(\sqrt N)}+\frac1N
\right),
\]
where \(\delta_{ij}\) is the Kronecker delta. The admissible values of \(N\), the range of \(n\), and the coefficient index ranges are not supplied. The assertion that these prescriptions define a state is part of the supplied formulation; no additional existence or positivity claim is established here.

The transition operator is
\[
P_\varphi=(\operatorname{id}\otimes\varphi)\circ\Delta.
\]
For a linear map \(P:A\to A\) on an operator space \(A\), its completely bounded norm is
\[
\|P\|_{\mathrm{cb}}
=\sup_{r\geq1}\bigl\|\operatorname{id}_{M_r}\otimes P:
M_r(A)\longrightarrow M_r(A)\bigr\|.
\]
The operator space or algebra completion on which this norm is intended is not specified in the supplied record.

\textbf{Source conjecture, with unresolved formulation.} Prove or disprove that the completely bounded norm of \(P_\varphi\) gives the correct cutoff parameter for the prescribed quantum random walks on \(S_N^+\). The original supplied wording refers more broadly to the relevant quantum-group random walks; the corrected supplied wording singles out the displayed family on \(S_N^+\).

The record supplies no equation relating the norm to a proposed cutoff time, no asymptotic range of \(N\), no distance to equilibrium, no initial condition, and no definition of cutoff in this setting. It also does not specify whether the intended norm involves the displayed operator itself, its iterates, or an operator measuring deviation from equilibrium. Consequently, the displayed prescription and the qualitative conjecture can be retained, but a uniquely quantified cutoff assertion cannot be recovered from this record alone.
\subsection{Short English statement}
The supplied problem asks whether the completely bounded transition norm determines cutoff for a specified family of random walks on \(S_N^+\). The state formula is retained exactly. The admissible parameters, representation indexing, norm setting, and meaning of the proposed cutoff parameter require clarification before the conjecture has a precise testable formulation.
\subsection{Sources}
\begin{itemize}
\item Local UnsolvedMath registry: supplied source record for permanent ID \(30003806\), titled \emph{Cutoff Parameters for Quantum-Group Random Walks}, in the Oberwolfach Reports Open Problems collection. Exact source problem code: \texttt{catalogue entry}. No proposer is named; the supplied proposed year is 2018.
\item Supplied report citation: \emph{Interactions between Operator Space Theory and Quantum Probability with Applications to Quantum Information} (2019), Oberwolfach Reports, pp.\ 1293--1379, DOI \texttt{10.4171/owr/2018/22}. \url{https://doi.org/10.4171/owr/2018/22}. The supplied background describes these reports as workshop records including open-problem discussions. No problem-specific page is supplied.
\item Supplied literature source, labeled ``Freslon on cutoff'': \url{https://arxiv.org/abs/1711.06555}.
\item Supplied literature source, labeled ``Later quantum group cutoff results'': \url{https://arxiv.org/abs/2010.03273}.
\item Literature assessment date: \texttt{2026-08-21}. This is supplied dated curation, not independently verified here. The supplied status is open, and the supplied assessment says that later cutoff results for free quantum groups do not prove the extracted conjecture that the completely bounded transition norm gives the correct cutoff parameter. The background explicitly requests independent verification of time-sensitive status and cited claims.
\item The registry labels the statement \texttt{corrected\_verified} and reports that it was checked directly against the official report, restoring the quantum permutation group, random-walk state, polynomial recurrence, and cutoff conjecture. This is a supplied verification claim, not an independent source check performed for this draft.
\end{itemize}
\end{document}
